\chapter{Sandbox（沙箱隔离）}

\vspace{0.5em}
\noindent\textit{上一章我们了解了 Memory 如何让智能体拥有持久记忆。但智能体执行的 Shell 命令、文件操作等
可能带来安全风险。本章将介绍 Sandbox（沙箱）隔离机制——通过 Docker 容器将工具执行限制在受控环境中，
确保智能体操作的安全性。}
\vspace{0.5em}

\section{概述}

OpenClaw 可以\textbf{在 Docker 容器内运行工具}以减少影响范围。沙箱隔离是\textbf{可选的}，由配置控制。Gateway 网关本身保留在主机上，仅工具执行在容器中隔离运行。

配置路径：\texttt{agents.defaults.sandbox.*}

\section{什么会被沙箱隔离}

\begin{center}
\begin{tabular}{ll}
\toprule
\textbf{隔离} & \textbf{不隔离} \\
\midrule
工具执行（exec, read, write, edit 等） & Gateway 网关进程本身 \\
沙箱浏览器（可选） & 提权工具（\texttt{tools.elevated}） \\
 & 明确允许在主机运行的工具 \\
\bottomrule
\end{tabular}
\end{center}

\section{沙箱模式}

\texttt{agents.defaults.sandbox.mode}：

\begin{center}
\begin{tabular}{ll}
\toprule
\textbf{模式} & \textbf{行为} \\
\midrule
\texttt{off} & 不使用沙箱 \\
\texttt{non-main} & 仅非主会话使用沙箱（默认） \\
\texttt{all} & 所有会话都在沙箱中运行 \\
\bottomrule
\end{tabular}
\end{center}

\section{作用域}

\texttt{agents.defaults.sandbox.scope}：

\begin{itemize}
  \item \texttt{session}（默认）：每个会话一个容器
  \item \texttt{agent}：每个智能体一个容器
  \item \texttt{shared}：所有沙箱会话共享一个容器
\end{itemize}

\section{工作区访问}

\texttt{agents.defaults.sandbox.workspaceAccess}：

\begin{center}
\begin{tabular}{ll}
\toprule
\textbf{值} & \textbf{行为} \\
\midrule
\texttt{none}（默认） & 使用沙箱工作区（\texttt{\textasciitilde{}/.openclaw/sandboxes}） \\
\texttt{ro} & 只读挂载智能体工作区 \\
\texttt{rw} & 读写挂载智能体工作区 \\
\bottomrule
\end{tabular}
\end{center}

\section{最小启用示例}

\begin{lstlisting}[language=json]
{
  agents: {
    defaults: {
      sandbox: {
        mode: "all",
        scope: "session",
        workspaceAccess: "none",
        docker: {
          image: "openclaw-sandbox:bookworm-slim"
        }
      }
    }
  }
}
\end{lstlisting}

\section{镜像构建}

\begin{lstlisting}[language=bash]
# Build sandbox image
scripts/sandbox-setup.sh

# Build sandbox browser image (optional)
scripts/sandbox-browser-setup.sh
\end{lstlisting}

默认沙箱容器\textbf{没有网络}，通过 \texttt{docker.network} 可覆盖。

\section{自定义绑定挂载}

将额外主机目录挂载到容器：

\begin{lstlisting}[language=json]
{
  agents: {
    defaults: {
      sandbox: {
        docker: {
          binds: [
            "/home/user/source:/source:ro",
            "/var/run/docker.sock:/var/run/docker.sock"
          ]
        }
      }
    }
  }
}
\end{lstlisting}

\section{setupCommand}

\texttt{setupCommand} 在容器创建后\textbf{运行一次}（非每次运行），用于安装额外依赖：

\begin{lstlisting}[language=json]
{
  agents: {
    defaults: {
      sandbox: {
        docker: {
          network: "bridge",
          setupCommand: "apt-get update && apt-get install -y python3"
        }
      }
    }
  }
}
\end{lstlisting}

\textbf{注意}：包安装需要网络出口（\texttt{network: "bridge"}）和可写根文件系统。

\section{多智能体覆盖}

每个智能体可独立覆盖沙箱配置：

\begin{lstlisting}[language=json]
{
  agents: {
    list: [{
      id: "build",
      sandbox: {
        mode: "all",
        docker: {
          binds: ["/mnt/cache:/cache:rw"]
        }
      }
    }]
  }
}
\end{lstlisting}

\section{调试}

\begin{lstlisting}[language=bash]
openclaw sandbox explain  # Check effective sandbox config
\end{lstlisting}

\section{本章小结}

本章介绍了 Sandbox 沙箱隔离的核心机制：

\begin{itemize}
  \item Sandbox 通过 \textbf{Docker 容器}将工具执行限制在受控环境中
  \item 支持 \textbf{off、non-main、all} 三种模式，及 \textbf{session、agent、shared} 三种作用域
  \item \textbf{Gateway 本身不被沙箱隔离}，仅工具执行在容器中运行
  \item 支持自定义挂载、初始化命令、多智能体独立覆盖
\end{itemize}

\noindent 至此，OpenClaw 的十大核心概念已经全部介绍完毕。下一章我们将进入\textbf{实战环节}——
在腾讯轻量云服务器上从零部署 OpenClaw，将全书所学知识付诸实践。
